\ifdefined\mainfile\else
\documentclass{article}
\usepackage{graphicx}
\usepackage{float}
\usepackage{amsmath}
\usepackage{amssymb}
\usepackage[a4paper, left=2.5cm, right=2.5cm, top=2.5cm, bottom=2.5cm]{geometry}
\usepackage[colorlinks=true,linkcolor=black,citecolor=blue,urlcolor=blue]{hyperref}
\begin{document}
\fi

\section{Communication technologies}
\label{sec:comms}

\subsection{WiFi (IEEE 802.11 b/g/n)}
\label{sec:wifi}

The bin connects to the home network over standard 2.4\,GHz WiFi, exposed by the ESP8266's built-in radio. WiFi uses CSMA/CA channel access on the unlicensed ISM band and provides per-link symmetric encryption (WPA2/WPA3). At link level the ESP8266 supports raw data rates from 1\,Mb/s up to 72.2\,Mb/s; in practice the limiting factor for our use case is the cloud round-trip, not the air link.

We chose WiFi (rather than e.g. BLE, LoRaWAN or Zigbee) for three reasons. First, the course explicitly requires it. Second, every household with a kitchen already has WiFi infrastructure, so the bin is plug-and-play. Third, WiFi is high-bandwidth and IP-routable, which means we can talk directly to any HTTP endpoint on the public internet (ThingSpeak) without needing a gateway.

\textbf{Pros.} Ubiquitous in homes; no extra gateway needed; high bandwidth; native IP stack so we can use ordinary HTTP REST calls.\\
\textbf{Cons.} Higher power consumption than BLE/LoRa, especially while the radio is active; depends on the user's home AP being reachable; 2.4\,GHz band can be noisy (microwaves, neighbours' APs).

\subsection{HTTP REST over the public internet}
\label{sec:http}

All cloud communication uses unencrypted HTTP GET requests against the ThingSpeak API. ThingSpeak is a simple cloud database with eight numbered fields per channel; writing a sample is a single URL fetch, and reading the most recent samples is another. Both the ESP and the dashboard speak the same protocol, so the integration code is symmetric.

\textbf{Pros.} Trivially simple; works from anything that can open a URL; ThingSpeak's free tier is sufficient for our cadence.\\
\textbf{Cons.} Plain HTTP, no encryption (acceptable for non-sensitive telemetry but not for production); each push is a full TCP+HTTP round-trip, which is heavy compared to a long-lived MQTT session.

\subsection{ThingSpeak React + IFTTT webhooks}
\label{sec:react}

ThingSpeak's \emph{React} feature watches a field for a condition (\texttt{field4 = 1} or \texttt{field2 < 3}) and fires a webhook URL when it triggers. The webhook target is an IFTTT applet, which in turn delivers a push notification to the IFTTT app on the user's phone.

\textbf{Pros.} No notification server of our own to maintain; native iOS/Android notifications; rules are configured in a browser without touching firmware.\\
\textbf{Cons.} Adds a hop (so notifications arrive with a few seconds of extra latency); free-tier React has a minimum 15-minute re-trigger interval; both services can change their APIs without warning.

\ifdefined\mainfile\else
\end{document}
\fi
